\section{Introduction}
\label{sec:introduction}

In a mean-payoff game, two players choose edges of a finite graph and
compete over the long-run average weight of the resulting infinite
walk. The objective is numerical, but the question of who can enforce
a nonnegative average is a finite decision problem. A central
algorithmic issue is the representation of the weights: a procedure
polynomial in their absolute values may take exponential time in the
length of the input. We give a deterministic algorithm whose running
time is quasipolynomial in that binary length and which determines
all winning starting vertices together.

\subsection{The problem and the main theorem}

Let $G=(V,E)$ be a finite directed graph with
$V=\{1,\ldots,n\}$ and $n\ge1$. Each vertex has at least one
outgoing edge and belongs to exactly one of
$V_{\mathrm{Max}}$ and $V_{\mathrm{Min}}$. The edge weights
$w:E\to\Z$ are arbitrary signed binary integers. Self-loops and
parallel edges are allowed; parallel edges are separate entries of
the explicit edge list. At a vertex its owner chooses the next edge.
Both players' strategies may depend on the entire finite history.
For an infinite play $\pi=(e_0,e_1,\ldots)$, write
\[
 \MP_w(\pi)=\liminf_{T\to\infty}\frac1T
                         \sum_{t=0}^{T-1}w(e_t).
\]
We count the complete input length $L$ in a fixed conventional binary
encoding of the graph, endpoints, ownership, and weights. For a start
$v$ and strategies $\sigma,\tau$, denote their induced play by
$\pi_{v,\sigma,\tau}$.

\begin{theorem}\label{thm:main}
A uniform deterministic algorithm returns exactly the set
\[
 U=\bigl\{v\in V:\exists\text{ Max strategy }\sigma\ 
       \forall\text{ Min strategies }\tau,
       \ \MP_w(\pi_{v,\sigma,\tau})\ge0\bigr\}
\]
for every game specified above. For an absolute constant $C>0$, the
complete computation uses at most
\[
                      2^{C(\log_2(L+2))^2}
\]
bit operations in a standard model polynomially equivalent to a
uniform deterministic Turing machine. The bound includes
preprocessing, exact arithmetic, recursion, and output of the entire
set $U$.
\end{theorem}

A positional strategy chooses its outgoing edge using only the current
vertex. Section~\ref{sec:full-solution} derives the exact rational
value at every starting vertex and globally optimal positional
strategies for both players, within the same quasipolynomial bit bound.
Section~\ref{sec:energy-credits} decides whether a prescribed initial
credit can keep every finite accumulated weight nonnegative in a
one-resource energy game, and computes minimum winning credits. Section~\ref{sec:tropical}
applies classical reductions to tropical feasibility and fixed-corank
tropical-rank tests. A polynomial bound for the complete binary input
would be a further improvement.

\subsection{Earlier algorithms and binary weights}

Ehrenfeucht and Mycielski proved the existence of positional optimal
strategies for mean-payoff games~\cite{EhrenfeuchtMycielski1979}.
Positional determinacy places the threshold problem in
$\mathrm{NP}\cap\mathrm{coNP}$. Zwick and Paterson developed
pseudopolynomial algorithms and complexity reductions
\cite{ZwickPaterson1996}. Such bounds
can be polynomial in the number of vertices and the largest absolute
weight while remaining exponential in the binary encoding length.
This distinction is central to Theorem~\ref{thm:main}.

The energy formulation asks whether a finite initial credit can keep
all finite accumulated weights nonnegative. Brim, Chaloupka, Doyen,
Gentilini, and Raskin developed bounded integer progress measures and
monotone lifting, with an $O(|E|nW_0)$ threshold and winning-region
bound when $W_0$ bounds the absolute weights
\cite[Theorem~8]{BrimEtAl2011}.
Dorfman, Kaplan, and Zwick developed a deterministic exponential-time
approach based on accelerated potential updates and scaling; Austin
and Dell'Erba corrected its update procedure
\cite{DorfmanKaplanZwick2019,AustinDellErba2023}. Kozachinskiy
gives a polyhedral interpretation and a weight-independent
$n^{O(1)}2^{n/2}$ arithmetic bound for energy winners
\cite{Kozachinskiy2021}.

Strategy-improvement algorithms provide another important comparison.
Bj\"orklund and Vorobyov proved an expected
$2^{O(\sqrt{n\log n})}$ arithmetic bound for the threshold partition
on graphs without parallel edges
\cite[Theorem~7.1]{BjorklundVorobyov2007}, following work with
Sandberg~\cite{BjorklundSandbergVorobyov2004}.
Recent analyses concern Random-Action-Removal~\cite{Zwick2026} and
Switch-All and Random-Edge for energy games
\cite{DorfmanKaplanZwick2026}. Iteration counts and arithmetic bounds
must be distinguished from bit bounds: strategy evaluation and the
sizes of intermediate numbers contribute to the latter.

Parity games use an extremal priority occurring infinitely often to
determine the winner. They reduce in polynomial time to mean-payoff
threshold games using weights of exponential magnitude but polynomial
binary length~\cite[Section~3]{Jurdzinski1998}. Calude, Jain,
Khoussainov, Li, and Stephan gave
a quasipolynomial algorithm for parity games
\cite{CaludeJainKhoussainovLiStephan2022}. The reduction goes from
parity to mean payoff, so this result does not by itself give a
quasipolynomial algorithm for arbitrary binary-weight mean-payoff games.
Daviaud, Jurdzi\'nski, and Lazi\'c combined succinct parity progress
measures with energy bounds in a pseudo-quasipolynomial algorithm for
mean-payoff parity games~\cite{DaviaudJurdzinskiLazic2018}; that
bound still contains a numerical weight parameter. Smoothed polynomial
results for independent Gaussian perturbations of payoffs on ergodic
game graphs concern a different input model~\cite{LoffSkomra2024}.

\subsection{How the proof works}

An integer potential assigns a number to each vertex. First replace
$w(e)$ by $\widehat w(e)=(n+1)w(e)+1$, which makes every simple
cycle sum nonzero while separating negative original sums from
nonnegative ones. For $z\in\Z^V$, let $F_i(z)$ be the maximum
of $\widehat w(e)+z_j$ over edges $e:i\to j$ at a Max vertex,
and the minimum at a Min vertex. For integer vectors $A\le u$,
the box $[A,u]$ consists of all integer vectors between them,
coordinatewise. In this box,
a subsolution is a vector $x\in[A,u]$ satisfying $x_i\le F_i(x)$
wherever $x_i>A_i$; a supersolution is a vector $y\in[A,u]$
satisfying $y_i\ge F_i(y)$ wherever $y_i<u_i$.
Clipping a vector to the box means replacing coordinate $z_i$ by
$\min\{u_i,\max\{A_i,z_i\}\}$. The comparison theorem in
\emph{Randomized quasipolynomial-time mean-payoff games}, our
companion article, gives a unique fixed point $h$ of $z\mapsto\min\{u,\max\{A,F(z)\}\}$,
and the inequalities $x\le h\le y$ for every such pair. These
finite-box facts and their basic iteration procedure are the only
results imported from the companion in the proof of
Theorem~\ref{thm:main}
\cite[Definition~2.2 and Lemmas~2.3--2.4]{OpenAI2026Randomized}.
Section~\ref{sec:boxes} states that interface in full.

The same section places each coordinate of $h$ near one of the
two boundaries of a sufficiently wide box $[0,D]^V$. The upper band
is winning for Max and the lower band for Min. Iterating the clipped
operator from the bottom can require a number of steps proportional
to the numerical width $D$. We instead ask which boundary labels are
forced, without computing $h$.

For a box $[A,A+D]$, give the vertices positive rational masses
$a_i$ and $b_i$. For a vertex set $S$, its masses are
$a(S)=\sum_{i\in S}a_i$ and $b(S)=\sum_{i\in S}b_i$. A Max witness is a subsolution whose support
$\{i:x_i>A_i\}$ has $a$-mass at most one; it requires a positive
label at coordinates near the upper boundary. A Min witness is the
dual supersolution with support $\{i:y_i<A_i+D\}$ of $b$-mass
at most one; it requires a negative label near the lower boundary.
The task is to meet every requirement of every witness at once. With
initial masses $1/n$, the box solution qualifies on both sides, so
these requirements determine every vertex's winner.

To solve that task recursively, we move an integer pivot vector $k$
near the middle of the box. For a Max witness, a common translation
aligns the largest deviation $x_i-k_i$ with the top of a narrower
box centered at $k$. Clipping at the bottom leaves only coordinates
whose deviations are close to that largest value. A dual operation
handles Min witnesses. Section~\ref{sec:translations} proves that
these operations preserve the required inequalities and restrict
supports to subsets of the original supports.

The algorithm in Section~\ref{sec:algorithm} first moves $k$
downward, then upward. Narrow recursive calls choose which
coordinates move. Each pass protects the extremal deviation of all
witnesses of one type. When a pass stops, any still-needed witness
of the other type must have a set of coordinates close to an extremum
with $a$-mass greater than $3/4$ for a Max witness or $b$-mass
greater than $3/4$ for a Min witness. Otherwise the last recursive call would have
forced another movement. This use of the final, stationary call is
the key deterministic certificate.

One call on a half-width box labels those large-mass sets. Its labels
need not settle the original boundary requirements. They instead
specify new masses: coordinates labelled in a witness's favor become
lighter, and the others heavier. Section~\ref{sec:correctness}
shows that every unsettled witness still has support mass at most
one. Meanwhile, every product $a_i b_i$ increases by a fixed factor.
Once every product exceeds one, each coordinate is excluded from
at least one type of witness support, and the labels can be chosen
directly. Starting with masses $1/n$, only $O(\log(n+1))$ such
product increases can therefore occur along a branch. The final call
on the original box uses the reweighted masses to settle the remaining
requirements. Every child except the half-width call increases the
products; that exceptional child instead reduces the width.
Section~\ref{sec:complexity} counts the possible positions of the
product increases along a branch and bounds the exact arithmetic in
each call, giving the stated bit complexity.

\subsection{Potential methods and recursive precision}

Gurvich, Karzanov, and Khachiyan developed potential transformations
in their algorithm for cyclic games
\cite{GurvichKarzanovKhachiyan1988,KarzanovLebedev1993}: adding
differences of endpoint potentials to edge weights preserves cycle
sums while changing the local inequalities. Paired potentials appear
in the work of Lifshits and Pavlov~\cite{LifshitsPavlov2007}.
Cadilhac, Casares, and Ohlmann give a general framework for accelerated
value iteration~\cite{CadilhacCasaresOhlmann2025}. These methods and
the progress-measure approach described above provide the setting
for our finite-box inequalities.

The closest recursive ancestry is Parys's reduced-precision parity
algorithm~\cite[Algorithm~2, Section~5, and Lemma~6.1]{Parys2019}
and its development by Lehtinen, Parys, Schewe, and Wojtczak
\cite{LehtinenParysScheweWojtczak2022}. Their guarantees protect
all dominions---sets within which one player can keep play and win---
within specified size bounds, and their recursion
separates one child retaining full precision from children of reduced
precision. We likewise use simultaneous obligations for small objects,
but smallness here means rational mass of a potential's support.
Translation near an extremal deviation, the stationary-pivot
certificates, and reweighting establish the needed progress for that
notion of smallness.

Ohlmann's symmetric mean-payoff recursion supplies closely related
potential-based context~\cite{Ohlmann2026}; its subexponential time
analysis is left open in the cited version. Jurdzi\'nski, Morvan,
Ohlmann, and Thejaswini use two simultaneous labellings for symmetric
acceleration of parity algorithms
\cite{JurdzinskiMorvanOhlmannThejaswini2020}; a related development
by Thejaswini, Ohlmann, and Jurdzi\'nski retains both decompositions
to accelerate symmetric recursion~\cite{ThejaswiniOhlmannJurdzinski2022}.
Arnold, Niwi\'nski, and Parys study recursive interval restriction in a
quasipolynomial black-box algorithm for nested fixed-point evaluation
\cite{ArnoldNiwinskiParys2021}. Their bound is expressed in
Boolean-lattice dimension and fixed-point nesting depth. The integer
box here has chains of length $nD+1$, so that result does not directly
give the dependence on $\log D$ established by our translation and
mass arguments.

Fijalkow, Gawrychowski, and Ohlmann studied mean-payoff value iteration
using universal graphs and lower bounds for those representations
\cite{FijalkowGawrychowskiOhlmann2020}. Those lower bounds constrain
universal-graph representations. The present recursion constructs and
reweights boxes for the given instance without using such a representation.

The companion algorithm~\cite{OpenAI2026Randomized} returns potential
vectors with probabilistic comparison guarantees stated separately
for each fixed admissible vector. Here the recursive guarantee is
simultaneous over all qualifying witnesses and concerns their boundary
labels. The two algorithms share the finite-box theory stated in
Section~\ref{sec:boxes}; the companion's randomized proof is
independent of the deterministic recursion. Section~\ref{sec:full-solution}
also explains how its winning-set interface yields randomized exact
values and positional strategies.
